\documentclass[11pt]{article}
\usepackage[a4paper,margin=1in]{geometry}
\usepackage{hyperref}
\usepackage{amsmath}
\usepackage{amssymb}
\usepackage{booktabs}
\usepackage{listings}
\usepackage{xcolor}

\definecolor{codebg}{RGB}{248,248,248}
\lstset{
  basicstyle=\ttfamily\small,
  backgroundcolor=\color{codebg},
  frame=single,
  breaklines=true,
  columns=fullflexible
}

\title{Three-Stage Legal Dialogue Training Pipeline Report\\(SFT $\rightarrow$ Multi-Objective $\rightarrow$ DPO)}
\author{Project: \texttt{q1\_3stage\_pipeline}}
\date{\today}

\begin{document}
\maketitle

\section{Executive Summary}
This report documents a strict \textbf{three-stage training pipeline} for a legal conversational model (Indian law: POCSO, IPC, CrPC) on \textbf{LLaMA-3.1-8B Instruct}. The stages are:
\begin{enumerate}
  \item \textbf{Stage 1 (SFT)}: supervised fine-tuning on multi-turn dialogue pairs (teacher-forced causal LM).
  \item \textbf{Stage 2 (Multi-objective)}: continue training from Stage 1 with a weighted combination of:
  \begin{itemize}
    \item $L_{\text{gen}}$ (LM loss),
    \item $L_{\text{entail}}$ (NLI teacher KL consistency),
    \item $L_{\text{triplet}}$ (embedding-space separation with dynamic hard negatives).
  \end{itemize}
  \item \textbf{Stage 3 (DPO)}: preference optimization using a reference model and dynamic rejected responses.
\end{enumerate}

We additionally implemented \textbf{reproducibility, checkpointed resumption, dynamic negative mining, speed optimizations for entailment}, and \textbf{optimizer stability diagnostics} (AMP unscale + clipping, gradient sanity stats, update ratio proxy, contribution ratios).

\section{Dataset and Splits}
\subsection{Data characteristics}
The dataset consists of multi-turn legal dialogues with Hindi/English code-mixing and Indian law domain constraints (statutes, procedures, and safety).

\subsection{Split protocol}
Training is dialogue-level split to avoid leakage across turns:
\begin{itemize}
  \item \textbf{Train}: \texttt{q1\_3stage\_pipeline/data/train\_70\_dialogues.jsonl}
  \item \textbf{Val}: \texttt{q1\_3stage\_pipeline/data/val\_10\_dialogues.jsonl}
  \item \textbf{Test}: \texttt{q1\_3stage\_pipeline/data/test\_20\_dialogues.jsonl}
\end{itemize}
Each dialogue JSONL record contains \texttt{turns}. Training code flattens dialogues into (prompt, response) examples in-memory.

\section{Model and Tokenization}
\subsection{Base model}
Default base model (from config):
\begin{lstlisting}
meta-llama/Meta-Llama-3.1-8B-Instruct
\end{lstlisting}

\subsection{Prompt format}
Prompts use a strict user/assistant delimiter format:
\begin{lstlisting}
[USER]: ...
[ASSISTANT]:
\end{lstlisting}
The assistant span is the supervised target; prompt tokens are masked in labels (\texttt{-100}).

\subsection{Code-mixed language tag (optional)}
To improve conditioning in code-mixed legal text, Stage 2 supports an optional prefix:
\begin{lstlisting}
--lang-tag-prefix "[HI_EN_LEGAL]"
\end{lstlisting}
This prepends the tag to every input prompt (train and val) without changing dataset files.

\section{Stage 1: Supervised Fine-Tuning (SFT)}
\subsection{Objective}
Standard teacher-forced causal LM cross-entropy:
\[
L_{\text{gen}} = \text{CE}(\text{LM}(x), y)
\]
where $x$ is the prompt and $y$ is the assistant completion. Prompt tokens are masked out in the label tensor.

\subsection{Artifacts}
Stage 1 outputs a Hugging Face-style \texttt{final/} directory that Stage 2 uses for initialization.

\section{Stage 2: Multi-Objective Training}
\subsection{High-level objective}
Stage 2 optimizes a weighted sum:
\[
L = L_{\text{gen}} + \lambda_e \, w_e(t)\,L_{\text{entail}} + \lambda_t \, w_t(t)\,L_{\text{triplet}}
\]

\paragraph{Dynamic weight scheduling.}
To avoid early instability from entailment spikes, the schedule uses step-progress $t \in [0,1]$:
\begin{itemize}
  \item Early (first 30\%): $w_e=0.3,\; w_t=0.3$
  \item Mid: $w_e=0.7,\; w_t=0.5$
  \item Late: $w_e=1.0,\; w_t=0.7$
\end{itemize}

\subsection{$L_{\text{entail}}$: NLI teacher KL}
We use a frozen DeBERTa MNLI teacher to produce a 3-way distribution (contradiction, neutral, entailment) for:
\[
(\text{premise}=y^+,\; \text{hypothesis}=\hat{y})
\]
where $y^+$ is the reference answer and $\hat{y}$ is the model's greedy completion for the prompt.

The student head predicts logits from pooled LM hidden states; the loss is:
\[
L_{\text{entail}} = \text{KL}\big(p_{\text{teacher}} \,\|\, p_{\text{student}}\big)
\]

\paragraph{Cost reductions.}
To reduce training time:
\begin{itemize}
  \item \textbf{Shorter greedy decode}: default \texttt{--entail-max-new-tokens 48} (recommended: 32--48).
  \item \textbf{Sparse computation}: \texttt{--entail-every N} computes entailment only every $N$ optimizer steps (default: 2).
  \item \textbf{LRU caching}: \texttt{--entail-cache-size} caches greedy $\hat{y}$ per prompt to avoid recomputation.
  \item \textbf{Teacher device}: \texttt{--nli-on-cpu} runs DeBERTa on CPU to save VRAM (slower but safer on shared GPUs).
\end{itemize}

\subsection{$L_{\text{triplet}}$: dynamic hard negatives + margin loss}
\paragraph{Dynamic negative generation.}
Negatives are \textbf{not stored in dataset files}. For each (prompt $x$, positive $y^+$) we form candidates:
\begin{itemize}
  \item \textbf{Model negative}: sampled generation from the current policy.
  \item \textbf{Rule corruption}: statute/procedure contradiction (e.g., IPC section swaps).
  \item \textbf{Cross-sample}: response from a different dialogue.
\end{itemize}

\paragraph{Hard / semi-hard mining.}
We embed candidates with a frozen sentence encoder (SBERT/MPNet) and select a hard negative that is:
\begin{itemize}
  \item sufficiently similar to the prompt (hard),
  \item less similar than the positive,
  \item optionally constrained to a \textbf{semi-hard band} relative to $\text{sim}(x,y^+)$.
\end{itemize}

\paragraph{Loss.}
Anchor is a learned projection of pooled LM hidden state; positives and negatives are sentence embeddings. With normalized vectors:
\[
L_{\text{triplet}} = \mathbb{E}\left[\max(0, d(a,p)-d(a,n)+m)\right]
\]
Default margin was reduced to $m=0.2$ (config can override) to reduce saturation.

\subsection{Checkpointing and resume}
Stage 2 writes a resumable checkpoint:
\begin{lstlisting}
output_dir/checkpoints/latest.pt
\end{lstlisting}
We optimized checkpoint size for large models by saving \textbf{PEFT adapter weights} (when applicable) rather than full 8B parameters.

\subsection{Training stability (AMP + gradients)}
\paragraph{Correct grad norm under AMP.}
With \texttt{GradScaler}, gradients are scaled. We therefore:
\begin{itemize}
  \item call \texttt{scaler.unscale\_(optimizer)} before measuring/clipping,
  \item clip gradients using \texttt{--grad-clip-max-norm} (default 1.0),
  \item log gradient sanity stats: \texttt{min\_grad}, \texttt{max\_grad}, \texttt{grad\_norm}.
\end{itemize}

\paragraph{Optional skip threshold.}
The ``skip step'' behavior is \textbf{disabled by default}. If desired:
\begin{lstlisting}
--skip-grad-norm-threshold 300
\end{lstlisting}
skips optimizer updates when the (unscaled, clipped) norm exceeds the threshold.

\paragraph{Update ratio proxy.}
We log:
\[
\text{update\_ratio} \approx \frac{\text{lr}\cdot\lVert g\rVert}{\lVert \theta\rVert}
\]
for trainable parameters, providing a compact stability signal.

\paragraph{Contribution ratios.}
Logs include per-step contributions:
\begin{itemize}
  \item \texttt{loss\_contrib\_gen}
  \item \texttt{loss\_contrib\_entail} (already includes schedule and $\lambda_e$)
  \item \texttt{loss\_contrib\_triplet} (includes schedule and $\lambda_t$)
  \item \texttt{triplet\_contrib\_frac} = triplet contribution / total
\end{itemize}

\subsection{Logging and evaluation hooks}
Training writes JSONL logs:
\begin{lstlisting}
output_dir/train_log.jsonl
\end{lstlisting}
Key fields include:
\begin{itemize}
  \item step counters: \texttt{epoch}, \texttt{opt\_step}, \texttt{micro\_step}
  \item losses: \texttt{loss}, \texttt{loss\_gen}, \texttt{loss\_entail}, \texttt{loss\_triplet}
  \item stability: \texttt{grad\_norm}, \texttt{min\_grad}, \texttt{max\_grad}, \texttt{scaler\_scale}, \texttt{update\_ratio}
  \item rolling stats (window=50): moving averages and spike/saturation rates
\end{itemize}

\paragraph{Fixed qualitative evaluation.}
Every \texttt{--fixed-eval-every} optimizer steps (default 200), we run 5 fixed prompts and log:
\begin{itemize}
  \item generated output
  \item teacher entailment probability (entailment label probability)
\end{itemize}
This provides a stable view of reasoning improvement during training.

\subsection{Debug-fast mode}
\texttt{--debug-fast} runs a cheap end-to-end sanity pass:
\begin{itemize}
  \item 100 training samples
  \item triplet disabled
  \item entailment computed every 5 steps
  \item generation caps reduced to 16 tokens
\end{itemize}

\section{Stage 3: DPO Preference Optimization}
\subsection{Goal}
Stage 3 performs preference optimization using:
\begin{itemize}
  \item \textbf{Policy}: Stage 2 final model
  \item \textbf{Reference}: frozen copy of Stage 2 (or saved reference)
  \item \textbf{Chosen}: ground truth $y^+$
  \item \textbf{Rejected}: dynamic hard negatives $y^-$ (same family as Stage 2)
\end{itemize}

\subsection{Beta sweep}
We support sweeping $\beta \in \{0.1, 0.5, 1.0\}$ (config) to tune the preference strength.

\section{Reproducibility and Project Layout}
\subsection{Reproducibility}
We set deterministic seeds and save run configuration files:
\begin{itemize}
  \item \texttt{config\_used.yaml}
  \item \texttt{run\_args.json}
\end{itemize}
for each training run directory under \texttt{q1\_3stage\_pipeline/logs/checkpoints/}.

\subsection{Key directories}
\begin{itemize}
  \item \texttt{q1\_3stage\_pipeline/configs/}: YAML configuration
  \item \texttt{q1\_3stage\_pipeline/data/}: JSONL splits (dialogue-level)
  \item \texttt{q1\_3stage\_pipeline/logs/checkpoints/stage1/}: Stage 1 checkpoints
  \item \texttt{q1\_3stage\_pipeline/logs/checkpoints/stage2/}: Stage 2 checkpoints + logs
  \item \texttt{q1\_3stage\_pipeline/logs/checkpoints/stage3/}: Stage 3 outputs (DPO)
\end{itemize}

\section{Operational Notes}
\subsection{Long epoch duration}
Stage 2 is slow because it includes:
\begin{itemize}
  \item on-the-fly generation for negatives,
  \item greedy decoding for entailment supervision,
  \item NLI teacher forward pass,
  \item embedding encoder for triplet mining.
\end{itemize}
The \texttt{--entail-every}, caching, and shorter \texttt{--entail-max-new-tokens} significantly reduce cost.

\subsection{Gated model warnings}
If the environment lacks Hugging Face access to gated model metadata, warnings about fetching remote config files may appear. Training can still proceed if model weights are already available locally.

\section{Recommended Defaults (Stage 2)}
\begin{lstlisting}
--entail-max-new-tokens 32 or 48
--entail-every 2 (or 3 if still slow)
--entail-cache-size 2048
--grad-clip-max-norm 1.0
--skip-grad-norm-threshold 0   (disabled)
\end{lstlisting}

\end{document}

